\documentclass[12pt,a4paper,oneside]{article}
\usepackage[brazil]{babel}
\usepackage[utf8]{inputenc}
\usepackage[dvipsnames]{xcolor}
\usepackage{listings}
\usepackage{wrapfig}
\usepackage{olymp}

\setcounter{problem}{1}

\begin{document}

\begin{problem}{Aceleração}{acelera}{2}

\begingroup
\setlength{\intextsep}{0pt}%
%\setlength{\columnsep}{0pt}

Em projetos de extensão da UFV, alunos de Engenharia Mecânica e interessados de outros cursos desenvolvem veículos e participam de competições nacionais.
A equipe Fórmula UFVolts Majorados projeta e fabrica veículos elétricos para participar do Fórmula SAE Brasil.
Já a equipe UFVbaja Pererecas projeta e constrói protótipos de veículos off-road para participar do Baja SAE Brasil.

Em algumas etapas, devem medir a aceleração do veículo. Se considerada uma aceleração constante, pode ser usada a Equação de Torricelli, proposta pelo físico e matemático italiano Evangelista Torricelli:

$$v^2 = v_0^2 + 2a\Delta x$$

onde $v$ representa a velocidade final do veículo e $v_0$ sua velocidade inicial, $a$ representa a aceleração e $\Delta x$ o deslocamento. Neste exercício, consideramos que o veículo parte do repouso, ou seja, $v_0 = 0$.

Como exemplo, estimando-se que o DeLorean do filme \textit{De Volta para o Futuro} tenha percorrido 370$m$ até atingir 88 milhas/hora (cerca de 142 km/h, equivalente a 39.444 m/s),
sua aceleração foi de $2.1$ m/s$^2$.

% \Notes

\InputFile

A entrada é composta por dois valores reais, $D$ e $V$, representando respectivamente o deslocamento do veículo em metros e sua velocidade final em km/h.
Ou seja, partindo do repouso ($v_0=0$), após percorrer $D$ metros sua velocidade é $V$ km/h.

Restrições: $1 \leq D \leq 1000$ e $1 \leq V \leq 200$.

\endgroup

\OutputFile

Escreva o valor da aceleração do veículo, em m/s$^2$, com uma casa decimal.

% \Notes

\Example

\begin{example}
\exmp{
100 36
}{
0.5
}%
\end{example}

\begin{example}
\exmp{
370 142
}{
2.1
}%
\end{example}

\begin{example}
\exmp{
12.3 45.6
}{
6.5
}%
\end{example}

%Comentários sobre o exemplo, se necessário.

\end{problem}

\end{document}